\begin{problem}[Warmup]{(a)}
    Suppose that $\alpha$ is the element of $S_7$ specified by \[\begin{pmatrix} 1 & 2 & 3 & 4 & 5 & 6 & 7 \\ 3 & 7 & 2 & 4 & 1 & 6 & 5 \end{pmatrix}.\] What is $\alpha(5)$? What is $\alpha(4)$? What is the inverse of $\alpha$?
\end{problem}
\unsolved

\begin{problem}[Warmup]{(b)}
    Suppose in addition that $\beta$ is the element of $S_7$ specified by \[\begin{pmatrix} 1 & 2 & 3 & 4 & 5 & 6 & 7 \\ 4 & 6 & 3 & 1 & 5 & 7 & 2 \end{pmatrix}.\] Calculate \[\beta^2, \quad \beta\alpha, \quad \alpha\beta, \quad \beta\alpha^{-1}, \quad \alpha\beta\alpha.\]
\end{problem}
\unsolved

\begin{problem}[Warmup]{(c)}
    Using the elements $\alpha$ and $\beta$ from the previous exercises, calculate ${(\alpha\beta)}^{-1}$ twice, once by using your computation of $\alpha\beta$ in Exercise b, and once by computing $\alpha^{-1}$ and $\beta^{-1}$, and then composing them (in the proper order!).
\end{problem}
\unsolved

\begin{problem}[Warmup]{(d)}
    How many elements are there in $S_4$? $S_5$? $S_6$?
\end{problem}
\unsolved

\begin{problem}[Warmup]{(e)}
    When someone says that every finite group is a group of permutations, what does that mean?
\end{problem}
\unsolved

\begin{problem}[Warmup]{(f)}
    Explain how the group of symmetries of the cube can be described as isomorphic to a subgroup of $S_{12}$, $S_8$, $S_6$, and $S_4$, depending on which set of objects associated with the cube you consider permuted.
\end{problem}
\unsolved

\begin{problem}[Warmup]{(g)}
    Every finite abelian group is isomorphic to a subgroup of a non-abelian group. Why?
\end{problem}
\unsolved

\begin{problem}{1}
    What are the orders of the permutations
    $\alpha, \beta, \beta^2, \beta\alpha, \alpha\beta, \beta\alpha^{-1}$, and
    $\alpha\beta\alpha$ given above in Exercises a and b? (Here
    \[\alpha = \begin{pmatrix} 1 & 2 & 3 & 4 & 5 & 6 & 7 \\ 3 & 7 & 2 & 4 & 1 & 6 & 5 \end{pmatrix}
    \quad\text{and}\quad
    \beta = \begin{pmatrix} 1 & 2 & 3 & 4 & 5 & 6 & 7 \\ 4 & 6 & 3 & 1 & 5 & 7 & 2 \end{pmatrix}
    \]
    in $S_7$.)
\end{problem}
\unsolved

\begin{problem}{2}
    Suppose that $m$ and $n$ are positive integers, and $m < n$. Define
    $I : S_m \to S_n$ as follows: Given $\alpha \in S_m$, we let $I(\alpha)(k) = \alpha(k)$,
    if $k \le m$, and $I(\alpha)(k) = k$, if $n \ge k > m$. Show that $I$ is a one-to-one
    homomorphism, which is not onto. \emph{Note:} You first must check that
    $I(\alpha) \in S_n$. We can paraphrase the contents of this exercise by asserting that
    (up to isomorphism) $S_m$ is a subgroup of $S_n$.
\end{problem}
\unsolved

\begin{problem}{3}
    Following the proof of Cayley's Theorem 29.1, determine explicitly which permutations
    of $S_4$ each of the elements of the group $\{1, -1, i, -i\}$ correspond to.
\end{problem}
\unsolved

\begin{problem}{4}
    Repeat Exercise 3 for the group of quaternions.
\end{problem}
\unsolved

\begin{problem}{5}
    Repeat Exercise 3 for $S_3$. Note that this gives us two representations of $S_3$ as a
    group of permutations: the original definition as a permutation of three elements, and
    the new one as a permutation of six elements!
\end{problem}
\unsolved

\begin{problem}{6}
    Let $n$ be a positive integer and let $k$ be a fixed integer, $1 \le k \le n$. Let
    \[G_k = \{\alpha \in S_n : \alpha(k) = k\}.\]
    Prove that $G_k$ is a subgroup of $S_n$. It is called a \textbf{stabilizer subgroup}
    of $S_n$. How many elements belong to $G_k$?
\end{problem}
\unsolved

\begin{problem}{7}
    Generalize Exercise 6: Let $K$ be any subset of $\{1, 2, \ldots, n\}$. Let
    \[G_K = \{\alpha \in S_n : \alpha(k) \in K, \text{ for all } k \in K\}.\]
    Prove that $G_K$ is a subgroup of $S_n$. How many elements belong to $G_K$?
\end{problem}
\unsolved
